\documentclass[tikz,border=6pt]{standalone}
\usepackage{ctex}
\usepackage{amsmath}
\usepackage{pgfplots}
\pgfplotsset{compat=1.18}
\usetikzlibrary{calc,angles,quotes,arrows.meta,positioning,decorations.markings}
\begin{document}
\begin{tikzpicture}[scale=1.0]
  % 观测点 O，水平距离 d=6，仰角 alpha=30度，高度 h=d*tan(30)
  \pgfmathsetmacro{\d}{6}
  \pgfmathsetmacro{\alphaang}{30}
  \pgfmathsetmacro{\h}{\d*tan(\alphaang)}

  \coordinate (O) at (0,0);          % 观测点
  \coordinate (B) at (\d,0);          % 塔底（与O同水平线）
  \coordinate (T) at (\d,\h);         % 塔顶（目标）

  % 地面 / 水平线（向右延长一点）
  \draw[->,thick] (-0.8,0) -- (\d+1.6,0) node[above right] {水平线};

  % 塔（竖直）
  \draw[very thick] (B) -- (T);
  % 塔身加一点阴影示意
  \fill[gray!18] ($(B)+(-0.12,0)$) rectangle ($(T)+(0.12,0)$);
  \draw[very thick] (B) -- (T);

  % 视线 O -> T
  \draw[thick] (O) -- (T);

  % 直角标记（塔底处）
  \draw (\d-0.32,0) -- (\d-0.32,0.32) -- (\d,0.32);

  % 仰角 alpha
  \pic[draw,thick,"$\alpha$",angle radius=1.25cm,angle eccentricity=1.35]
       {angle = B--O--T};

  % 水平距离 d 标注（在地面下方）
  \draw[<->,thin] (0,-0.55) -- (\d,-0.55)
        node[midway,fill=white,inner sep=1pt] {水平距离 $d$};

  % 高度 h 标注（塔右侧）
  \draw[<->,thin] (\d+0.55,0) -- (\d+0.55,\h)
        node[midway,fill=white,inner sep=1pt] {$h=d\tan\alpha$};

  % 点标注
  \fill (O) circle (1.6pt) node[left=2pt] {观测点 $O$};
  \fill (B) circle (1.6pt) node[below right=1pt and 1pt] {塔底 $B$};
  \fill (T) circle (1.6pt) node[above] {塔顶 $T$};

  % 半透明说明框
  \node[fill=yellow!18,fill opacity=0.75,text opacity=1,rounded corners,
        inner sep=4pt,align=left] at (2.05,2.55)
       {仰角 $\alpha$：视线与\\水平线的夹角};
\end{tikzpicture}
\end{document}
